%%%%%%
%$Autor: Wings$
%$Projekt: A25-12$
%$Dateiname: ArduinoIDE2Mac.tex$
%$Beschreibung: Installation und Überblick zu Funktionen und Tools der Arduino IDE fir OS Geräte$
%%%%%%


\subsection{Installation der Arduino IDE für MacOS}

Der Download der Arduino IDE erfolgt über die Website \HREF{Arduino.cc}. Geben Sie dazu 
\glqq Arduino IDE Download \grqq in Ihren Browser ein (vgl. Abbildung~\ref{pic:IDEMac01}).

\begin{center}
	\includegraphics[width=1\textwidth,keepaspectratio]{Arduino/ArdiunoIDE2/Mac/IDEMac01}
	\captionof{figure}{Aufruf der Arduino-Webseite über die Suchmaschine}
	\label{pic:IDEMac01}
\end{center}

Abbildung~\ref{pic:IDEMac02} zeigt die Ergebnisse der Websuche. Klicken Sie auf den Eintrag mit \glqq Arduino.cc\grqq im Adresskopf.

\begin{center}
	\includegraphics[width=1\textwidth,keepaspectratio]{Arduino/ArdiunoIDE2/Mac/IDEMac02}
	\captionof{figure}{Ergebnisse der Websuche zur Arduino IDE}
	\label{pic:IDEMac02}
\end{center}

Wählen Sie die passende Version für Ihren Mac aus (vgl. Abbildung~\ref{pic:IDEMac03}).

\begin{center}
	\includegraphics[width=1\textwidth,keepaspectratio]{Arduino/ArdiunoIDE2/Mac/IDEMac03}
	\captionof{figure}{Download der Arduino IDE starten}
	\label{pic:IDEMac03}
\end{center}

Es erscheint eine Spendenanfrage. Sie können freiwillig spenden oder über \glqq JUST DOWNLOAD \grqq fortfahren (vgl. Abbildung~\ref{pic:IDEMac04}).

\begin{center}
	\includegraphics[width=1\textwidth,keepaspectratio]{Arduino/ArdiunoIDE2/Mac/IDEMac04}
	\captionof{figure}{Optional: Spendenaufruf beim Download}
	\label{pic:IDEMac04}
\end{center}

Vor dem Download können Sie optional den Newsletter abonnieren, vgl. Abbildung~\ref{pic:IDEMac05}.

\begin{center}
	\includegraphics[width=1\textwidth,keepaspectratio]{Arduino/ArdiunoIDE2/Mac/IDEMac05}
	\captionof{figure}{Optional: Newsletter abonnieren}
	\label{pic:IDEMac05}
\end{center}

Da es sich um ein Programm aus dem Internet handelt, muss der Download zunächst erlaubt werden, vgl. Abbildung~\ref{pic:IDEMac06}.

\begin{center}
	\includegraphics[width=1\textwidth,keepaspectratio]{Arduino/ArdiunoIDE2/Mac/IDEMac06}
	\captionof{figure}{Download-Erlaubnis erteilen}
	\label{pic:IDEMac06}
\end{center}

Über den Pfeil oben rechts lässt sich der Fortschritt anzeigen (vgl. Abbildung~\ref{pic:IDEMac07}).

\begin{center}
	\includegraphics[width=1\textwidth,keepaspectratio]{Arduino/ArdiunoIDE2/Mac/IDEMac07}
	\captionof{figure}{Download-Fortschritt einsehen}
	\label{pic:IDEMac07}
\end{center}

Öffnen Sie nach dem Download die Datei per Doppelklick (vgl. Abbildung~\ref{pic:IDEMac08}).

\begin{center}
	\includegraphics[width=1\textwidth,keepaspectratio]{Arduino/ArdiunoIDE2/Mac/IDEMac08}
	\captionof{figure}{Heruntergeladene Datei öffnen}
	\label{pic:IDEMac08}
\end{center}

Ziehen Sie nun die App in den Programme-Ordner (vgl. Abbildung~\ref{pic:IDEMac10}).

\begin{center}
	\includegraphics[width=1\textwidth,keepaspectratio]{Arduino/ArdiunoIDE2/Mac/IDEMac10}
	\captionof{figure}{Installation durch Ziehen in den Programme-Ordner}
	\label{pic:IDEMac10}
\end{center}

Die Installation ist abgeschlossen. Öffnen Sie die App über den Programme-Ordner (vgl. Abbildung~\ref{pic:IDEMac11}).

\begin{center}
	\includegraphics[width=1\textwidth,keepaspectratio]{Arduino/ArdiunoIDE2/Mac/IDEMac11}
	\captionof{figure}{Arduino IDE starten}
	\label{pic:IDEMac11}
\end{center}

Bestätigen Sie die Sicherheitsabfrage mit \glqq Öffnen\grqq (vgl. Abbildung~\ref{pic:IDEMac12}).

\begin{center}
	\includegraphics[width=1\textwidth,keepaspectratio]{Arduino/ArdiunoIDE2/Mac/IDEMac12}
	\captionof{figure}{Sicherheitsabfrage bei erstmaligem Start}
	\label{pic:IDEMac12}
\end{center}

\subsection{Bibliotheken installieren}
\label{Bibliotheken installieren Mac}

Klicken Sie auf das Buch-Symbol (vgl. Abbildung~\ref{pic:IDEMac13}) und suchen Sie die gewünschte Bibliothek. Die vollständige Liste finden Sie in Kapitel \ref{Bibliotheken}. Installieren Sie jede Bibliothek durch Klick auf \glqq Install \grqq.

\begin{center}
	\includegraphics[width=1\textwidth,keepaspectratio]{Arduino/ArdiunoIDE2/Mac/IDEMac13}
	\captionof{figure}{Bibliotheken-Suchmaske der Arduino IDE}
	\label{pic:IDEMac13}
\end{center}

Abbildung~\ref{pic:IDEMac14} zeigt eine erfolgreich installierte Bibliothek.

\begin{center}
	\includegraphics[width=1\textwidth,keepaspectratio]{Arduino/ArdiunoIDE2/Mac/IDEMac14}
	\captionof{figure}{Erfolgreiche Installation einer Bibliothek}
	\label{pic:IDEMac14}
\end{center}

\subsection{Board Manager}

Mit dem Board Manager können Sie zusätzliche Boards hinzufügen (vgl. Abbildung~\ref{pic:IDEMac15}). Suchen Sie zum Beispiel nach dem \glqq ESP \grqq oder \glqq Nano 33 BLE Sense \grqq und installieren Sie das passende Board. 
\begin{center}
	\includegraphics[width=1\textwidth,keepaspectratio]{Arduino/ArdiunoIDE2/Mac/IDEMac15}
	\captionof{figure}{Der Board Manager in der Arduino IDE}
	\label{pic:IDEMac15}
\end{center}

\label{PortauswahlMAC}
Klappen Sie das USB-Menü aus (vgl. Abbildung~\ref{pic:IDEMac16}) und wählen Sie \glqq Select other board and port\grqq.

\begin{center}
	\includegraphics[width=1\textwidth,keepaspectratio]{Arduino/ArdiunoIDE2/Mac/IDEMac16}
	\captionof{figure}{Board- und Portauswahl öffnen}
	\label{pic:IDEMac16}
\end{center}

Geben Sie \glqq Nano 33 BLE Sense\grqq in das Suchfeld ein, verbinden Sie das Board und wählen Sie den richtigen Port (vgl. Abbildung~\ref{pic:IDEMac17}).

\begin{center}
	\includegraphics[width=1\textwidth,keepaspectratio]{Arduino/ArdiunoIDE2/Mac/IDEMac17}
	\captionof{figure}{Passenden Board und Port auswählen}
	\label{pic:IDEMac17}
\end{center}

\subsection{Beispielprogramm \glqq Blink\grqq}

Wählen Sie den Pfad gemäß Abbildung~\ref{pic:IDEMac18}, um ein Beispielprogramm zu laden.

\begin{center}
	\includegraphics[width=1\textwidth,keepaspectratio]{Arduino/ArdiunoIDE2/Mac/IDEMac18}
	\captionof{figure}{Beispielprogramme öffnen}
	\label{pic:IDEMac18}
\end{center}

Das Blink-Programm wird geöffnet (vgl. Abbildung~\ref{pic:IDEMac19}) und kann nach Auswahl des Ports hochgeladen werden (siehe \ref{PortauswahlMAC}). Die Kompilierung erfolgt automatisch vor dem Upload oder manuell über das Häkchen-Symbol.

\begin{center}
	\includegraphics[width=1\textwidth,keepaspectratio]{Arduino/ArdiunoIDE2/Mac/IDEMac19}
	\captionof{figure}{Das geöffnete Beispielprogramm \glqq Blink\grqq}
	\label{pic:IDEMac19}
\end{center}

Nach erfolgreichem Upload erscheint ein Terminalfenster mit einer Bestätigung.



